\subsection*{NT-A20--A21. Euclidean ideal classes and one-point localization}

\paragraph{NT-A20: the exact Picard-group criterion.}
Let $K$ be a number field, let $S$ be a finite set of places
containing every archimedean place, and put $R=\mathcal O_{K,S}$.
Assume $|S|\ge2$. For every positive squarefree integer $m$, set
\[
 E_m=K\bigl(\zeta_m,(R^\times)^{1/m}\bigr).
 \tag{A20.1}
\]
These are finite number fields: the finitely generated group
$R^\times$ has only finitely many generators whose $m$th roots
need to be adjoined. Assume precisely RH for all $\zeta_{E_m}$.
Input $M$ supplies this family through its trivial Hecke characters.

A fractional ideal class $C\in\operatorname{Pic}(R)$ is
Euclidean in Lenstra's sense if and only if
\[
                 \operatorname{Pic}(R)=\langle C\rangle.
 \tag{A20.2}
\]
Thus, if $\operatorname{Pic}(R)$ is cyclic of order $h_S$,
it has exactly $\varphi(h_S)$ Euclidean ideal classes,
where $\varphi(1)=1$; otherwise it has none.
At most one of these classes can be Euclidean for the ideal norm.
In particular, when $h_S>2$ is cyclic, at least
$\varphi(h_S)-1$ classes admit a Euclidean reduction but
do not admit the norm reduction.

Here is the precise definition and the full elementary part
of the argument. Write
\[
 \mathcal E(R)=\{I:\ I\text{ is a nonzero fractional ideal of }R,\
                                      R\subseteq I\}.
\]
For a representative $\mathfrak c$ of $C$, a Euclidean
reduction is a map $\psi:\mathcal E(R)\to W$ to a well-ordered
set such that for every $I\in\mathcal E(R)$ and
$x\in I\mathfrak c\setminus\mathfrak c$, there is
$y\in\mathfrak c$ satisfying
\[
 J=(x-y)^{-1}I\mathfrak c\in\mathcal E(R),
                    \qquad \psi(J)<\psi(I).
 \tag{A20.3}
\]
The element $x-y$ is nonzero. The displayed membership
holds because $x-y\in I\mathfrak c$.
Replacing $\mathfrak c$ by a principal multiple leaves
the definition unchanged.

\begin{proof}
For necessity, suppose there is an ideal class outside
$\langle C\rangle$. Every fractional ideal class has a
representative in $\mathcal E(R)$: for a nonzero
$a$ in a fractional ideal $\mathfrak a$, the ideal
$a^{-1}\mathfrak a$ contains $R$.
Choose $I\in\mathcal E(R)$ outside that subgroup with
$\psi(I)$ smallest. Then $I\ne R$, and invertibility
of $\mathfrak c$ gives $I\mathfrak c\ne\mathfrak c$.
Choose $x$ in their difference and apply (A20.3).
Since $[J]=[I]C$, the class $[J]$ is still outside
$\langle C\rangle$, contradicting the minimal choice.
This proves necessity without RH.

For sufficiency, apply Lenstra's Proposition~(2.2),
printed p.~127 of \emph{Euclidean ideal classes},
Ast\'erisque~61 (1979), 121--131. Its hypotheses are
exactly $|S|\ge2$ and RH for (A20.1), and its conclusion
is that every class generating $\operatorname{Pic}(R)$
has such a reduction. The printed inequality is
$|S|\ge2$; the searchable OCR incorrectly renders it
as a strict inequality. This is a cited classical
analytic input. The paper refers its sufficiency
argument to the methods of Lenstra~1977, reference~[9];
it does not print a full new analytic proof at p.~127.
We do not represent the statement check as a fresh
verification of that omitted analytic argument.

The number of generators of a finite cyclic group of
order $h_S$ is $\varphi(h_S)$, giving the count.
For the norm assertion, use the single fixed map
\[
                       \psi_N(I)=N_R(I)^{-1}\in\mathbf N.
 \tag{A20.4}
\]
It is integral because $I^{-1}$ is integral and
$N_R(I)^{-1}=\#(R/I^{-1})$.
If two classes were Euclidean for this same map,
choose an $I\ne R$ of smallest $\psi_N(I)$.
For either class, (A20.3) produces a $J$ with smaller
value. Minimality among nontrivial overideals forces
$J=R$, so both classes must equal $[I]^{-1}$.
The set from which $I$ is chosen is nonempty:
there is a prime ideal outside the finite set $S$.
Finally,
\[
 \psi_N\bigl((x-y)^{-1}I\mathfrak c\bigr)<\psi_N(I)
                  \quad\Longleftrightarrow\quad
                     N_R(x-y)<N_R(\mathfrak c),
\]
which is exactly norm-Euclidean reduction. This proves
the claimed uniqueness and the lower count of classes
that are not norm-Euclidean.
\end{proof}

The necessity argument and fixed-function uniqueness
are elementary reformulations of Lenstra's
Lemma~(1.5), Theorem~(1.6), and the paragraph after
Exercise~(1.7), printed pp.~124--125; those complete
elementary proofs were read. The mathematical content
of the sufficiency is classical. The explicit family,
class count, and comparison with $M$ here make its
input and consequences reviewable. No finite-order
Dirichlet input alone supplies the general unit
Kummer family. No complexity bound for evaluating
the well-ordered reduction is claimed.

\paragraph{NT-A21: a bounded single closed point makes a Euclidean ring.}
Let $K$ be any number field whose ordinary class group
is cyclic, of order $h$. Write $n=[K:\mathbf Q]$ and
$D=|\operatorname{disc}K|$. Under $M$, for every
generator $C$ of $\operatorname{Cl}(K)$ there is a
prime ideal $\mathfrak p$ of absolute residue degree
one such that
\[
 [\mathfrak p]=C,\qquad
 N\mathfrak p\le
                \bigl[h(4\log D+2.5n)+5\bigr]^2,
 \tag{A21.1}
\]
and the ring
\[
             R=\mathcal O_{K,S_\infty\cup\{\mathfrak p\}}
 \tag{A21.2}
\]
is Euclidean. In particular, this holds for fields
whose original unit group is finite. The norm bound
controls the one closed point removed; it is not a
bound for the running time of the Euclidean reduction.

The exact analytic inputs, weaker than universal $M$,
are RH for $\zeta_{\mathbf Q}$ and for the ordinary
Hilbert class-field zeta function $\zeta_{H_K}$,
followed by RH for the fields (A20.1) formed from
the particular ring (A21.2). The latter family
depends on the selected prime. No claim that the
first two zeta functions imply that family is made.

\begin{proof}
NT-C07, the checked specialization of Bach--Sorenson
Theorem~5.1, gives (A21.1). Its discriminant
substitution is $D_{H_K}=D^h$ and $[H_K:\mathbf Q]=hn$.
For $K=\mathbf Q$ take $\mathfrak p=(2)$ directly;
it satisfies (A21.1).

Extension of ideals to $R$ gives the exact quotient
\[
 \operatorname{Pic}(R)=
       \operatorname{Cl}(K)/\langle[\mathfrak p]\rangle.
 \tag{A21.3}
\]
Indeed every ideal of $R$ is represented by a finite
product of primes different from $\mathfrak p$,
so extension is onto. If an ideal becomes principal,
its divisor differs from a principal divisor only at
$\mathfrak p$, showing that the kernel is the displayed
cyclic subgroup. Conversely every power of
$\mathfrak p$ becomes the unit ideal.
Since $[\mathfrak p]$ generates $\operatorname{Cl}(K)$,
the group in (A21.3) is zero.

There are at least two places in
$S_\infty\cup\{\mathfrak p\}$. Apply NT-A20 to its
principal ideal class; equivalently apply the
classical PID statement NT-A03. This makes $R$
Euclidean. For clarity, the equivalence between a
principal Euclidean ideal class and a Euclidean domain
can also be seen directly. For nonzero $b\in R$ put
$\phi(b)=\psi(b^{-1}R)$, with $\psi$ from (A20.3) and
$\mathfrak c=R$. If $a/b\notin R$, apply (A20.3) to
$I=b^{-1}R$ and $x=a/b$ to find $y\in R$. Then
$r=a-by\ne0$ and
\[
       (a/b-y)^{-1}b^{-1}R=r^{-1}R,\qquad
                                  \phi(r)<\phi(b).
 \tag{A21.4}
\]
If $a/b\in R$, the remainder is zero, so the usual
division property holds in all cases.

Finally $\mathfrak p^h=(a)$ for some $a\in\mathcal O_K$,
because $[\mathfrak p]$ has order $h$. Its divisor
is supported only at $\mathfrak p$, and hence
$R=\mathcal O_K[1/a]$. Thus (A21.2) is the actual
affine open obtained by deleting this single closed
point, even when $\mathfrak p$ is not principal.
\end{proof}

Geometrically, NT-A20 concerns reduction in an
invertible sheaf generating
$\operatorname{Pic}(\operatorname{Spec}R)$.
NT-A21 removes one quantitatively bounded closed point
whose divisor class kills that Picard group, and then
uses the unit-Kummer RH family to supply Euclidean
reduction. Cyclicity remains a hypothesis; the
argument neither asserts that all class groups are
cyclic nor that the ideal norm is a Euclidean function.
